\section{Approximation of log-concave densities}\label{geo:section}

The smooth estimate of Theorem~\ref{mat:entropy-smooth} requires a
potential with positive, bounded Hessian. We remove those restrictions
while preserving the entropy, covariance, and first moments. The same
approximation will let us study equality by applying the strict
remainder estimate to smooth densities whose entropy gaps tend to zero.

\subsection{Convex approximation with entropy convergence}

We use extended-valued convex potentials: $V:\R^m\to(-\infty,+\infty]$ is
proper if it is finite somewhere and never takes the value $-\infty$.
The convention $e^{-\infty}=0$ allows bounded supports.
The elementary tail estimate below supplies a common integrable bound for
the approximation, including its entropy integrand.

\begin{lemma}\label{geo:coercivity}
Let $V:\R^m\to(-\infty,+\infty]$ be proper, lower semicontinuous, and convex,
with
\[
  0<\int_{\R^m}e^{-V(x)}\,\dd x<\infty.
\]
Then there are $a>0$ and $b\geq0$ such that
\begin{equation}\label{geo:linear-coercivity}
  V(x)\geq a|x|-b\qquad(x\in\R^m).
\end{equation}
In particular, $e^{-V}$ has moments of every order, and
$\int |V|e^{-V}<\infty$, where $|V|e^{-V}$ is interpreted as zero when
$V=+\infty$.
\end{lemma}

\begin{proof}
The convex set $\{V<\infty\}$ has positive measure and hence nonempty
interior: otherwise its affine hull would be a proper affine subspace.
Choose a simplex with all vertices in this set and with nonempty interior.
Convexity bounds $V$ above by the largest of its values at the vertices
throughout the simplex.  Consequently, some ball $B(x_0,r)$ lies in a
sublevel set $\{V\leq M\}$, where $M\geq V(x_0)$.

The closed convex set $A=\{V\leq M+1\}$ has finite volume, since
\[
  |A|\leq e^{M+1}\int e^{-V}.
\]
It is bounded.  Indeed, the convex hull of $B(x_0,r)$ and a point $z\in A$
contains a cone with base an $(m-1)$-dimensional disk of radius $r$ through
$x_0$, perpendicular to $z-x_0$, and height $|z-x_0|$.  Its volume tends to
infinity with $|z-x_0|$.  In dimension one the same assertion follows from
the length of the interval joining $x_0$ to $z$.

Choose $R>r$ so large that $A\subset B(x_0,R)$.  If
$d=|x-x_0|\geq R$ and $V(x)<\infty$, put
$y=x_0+(R/d)(x-x_0)$.  Then $y\notin A$, and convexity gives
\[
  M+1<V(y)\leq \left(1-\frac Rd\right)V(x_0)+\frac RdV(x).
\]
Thus $V(x)>V(x_0)+d/R$.  The inequality is automatic if $V(x)=+\infty$.
On the compact ball $\overline B(x_0,R)$, lower semicontinuity and properness
give a finite lower bound: a sequence with values tending to $-\infty$
would have a convergent subsequence and violate lower semicontinuity.
Combining these two bounds proves \eqref{geo:linear-coercivity}, after
enlarging $b$.

The moment assertion follows from $e^{-V(x)}\leq e^b e^{-a|x|}$.
For the entropy assertion, the elementary bound
\begin{equation}\label{geo:entropy-domination}
  |u|e^{-u}\leq C_b e^{-u/2}\qquad(u\geq-b)
\end{equation}
holds because $|u|e^{-u/2}$ is bounded on $[-b,\infty)$.
Substitute $u=V(x)$ and use \eqref{geo:linear-coercivity}.
\end{proof}

\begin{lemma}\label{geo:approximation}
Let $f=e^{-V}$ be a log-concave probability density on $\R^m$, represented
by a proper lower semicontinuous convex potential $V$.  There are smooth
convex potentials $V_j$ and probability densities
\[
  f_j=Z_j^{-1}e^{-V_j},\qquad Z_j=\int_{\R^m}e^{-V_j},
\]
such that
\begin{equation}\label{geo:hessian-bounds}
  \frac2j\Id\preceq D^2V_j\preceq\left(j+\frac2j\right)\Id
  \qquad(j\geq1),
\end{equation}
and
\begin{equation}\label{geo:approximation-limits}
  Z_j\longrightarrow1,\qquad \norm{f_j-f}_{L^1}\longrightarrow0,\qquad
  h(f_j)\longrightarrow h(f),\qquad
  \Cov(f_j)\longrightarrow\Cov(f).
\end{equation}
The covariance limit is entrywise, the means of $f_j$ converge to the mean
of $f$, and all these quantities are finite.
\end{lemma}

\begin{proof}
Fix $a,b$ from Lemma~\ref{geo:coercivity}.  For each positive integer $j$,
define the Moreau envelope \cite[Section~7]{Moreau1965}
\begin{equation}\label{geo:moreau}
  Q_j(x)=\inf_{y\in\R^m}\left\{V(y)+\frac j2|x-y|^2\right\}.
\end{equation}
The expression in braces is lower semicontinuous, strictly convex on its
effective domain, and tends to infinity as $|y|\to\infty$.  It therefore
has a unique minimizer $p_j(x)$, and $Q_j$ is finite everywhere.  Moreover,
\begin{equation}\label{geo:envelope-lower}
  Q_j(x)\geq a|x|-b-\frac{a^2}{2j}.
\end{equation}
Indeed, $V(y)\geq a|x|-a|x-y|-b$, and minimizing
$jr^2/2-ar$ over $r\geq0$ proves the bound.

For completeness, the needed regularity of $Q_j$ follows directly from
the minimization.  The optimality condition is
$j(x-p_j(x))\in\partial V(p_j(x))$, where $\partial V$ denotes the convex
subdifferential.  To obtain it, compare the minimum at $p=p_j(x)$ with
$p+s(z-p)$, use convexity to bound $V(p+s(z-p))$ above, divide by $s>0$,
and let $s\downarrow0$.  The resulting inequality is
$V(z)\geq V(p)+j\langle x-p,z-p\rangle$.
Adding the two subgradient inequalities at $p_j(x)$ and
$p_j(x')$ gives
\[
  \langle x-x',p_j(x)-p_j(x')\rangle
  \geq |p_j(x)-p_j(x')|^2.
\]
Consequently both $p_j$ and $x\mapsto x-p_j(x)$ are $1$-Lipschitz.
Evaluating the infimum at $p_j(x)$ and, in the reverse direction, at
$p_j(x+h)$ shows that
\[
  \nabla Q_j(x)=j(x-p_j(x)).
\]
For example, the first comparison gives an upper error $j|h|^2/2$
after subtracting $j\langle x-p_j(x),h\rangle$, and the second gives
a lower error of order $j|h|^2$ by the Lipschitz estimate just proved.
Thus $Q_j$ is continuously differentiable with $j$-Lipschitz gradient.
It is convex, since the function minimized in \eqref{geo:moreau} is jointly
convex in $(x,y)$.

We next check convergence before smoothing.  The functions $Q_j$ increase
with $j$, and $Q_j(x)\leq V(x)$ when $V(x)$ is finite.  If $Q_j(x)$ stays
bounded above as $j\to\infty$, the identity at the minimizer and $V\geq-b$
imply
\[
  |x-p_j(x)|^2\leq\frac{2(Q_j(x)+b)}j\longrightarrow0.
\]
Lower semicontinuity then gives
$V(x)\leq\liminf_j V(p_j(x))\leq\lim_jQ_j(x)$.
This proves
\begin{equation}\label{geo:envelope-limit}
  Q_j(x)\longrightarrow V(x)\quad\hbox{for every }x,
\end{equation}
including the value $+\infty$.

Choose a nonnegative even smooth function $\eta$ supported in the unit ball
with integral one, and set $\eta_\varepsilon(x)=\varepsilon^{-m}
\eta(x/\varepsilon)$.  Define
\begin{equation}\label{geo:smooth-potentials}
  V_j=Q_j*\eta_{1/j}+\frac{|x|^2}j.
\end{equation}
Convolution preserves convexity and the Lipschitz bound on the gradient;
hence \eqref{geo:hessian-bounds} holds.  Evenness gives
$\int z\eta(z)\,\dd z=0$, so Jensen's inequality and
\eqref{geo:envelope-lower} yield the common bound
\begin{equation}\label{geo:common-tail}
  V_j(x)\geq Q_j(x)\geq a|x|-B,
  \qquad B=b+\frac{a^2}2.
\end{equation}

Smoothing at this scale does not affect the limit.  The $j$-Lipschitz
bound on the gradient gives
\[
  Q_j(x-z)\leq Q_j(x)-\langle\nabla Q_j(x),z\rangle+\frac j2|z|^2.
\]
Integrating against $\eta_{1/j}$ cancels the linear term, so
\[
  0\leq (Q_j*\eta_{1/j})(x)-Q_j(x)\leq\frac1{2j}
  \qquad(x\in\R^m).
\]
Together with \eqref{geo:envelope-limit}, this proves
$V_j(x)\to V(x)$ everywhere.

Now \eqref{geo:common-tail} dominates $(1+|x|^2)e^{-V_j(x)}$ by an
integrable function independent of $j$.  Applying
\eqref{geo:entropy-domination} with $B$ in place of $b$ also gives
\[
  |V_j(x)|e^{-V_j(x)}\leq C e^{-a|x|/2}.
\]
At points where $V=+\infty$, both $e^{-V_j}$ and $V_je^{-V_j}$ tend to
zero.  Dominated convergence therefore proves convergence of the
normalizing constants, first moments, second moments, and the integrals
$\int V_je^{-V_j}$.  Finally,
\[
  h(f_j)=\frac1{Z_j}\int V_je^{-V_j}+\log Z_j,
\]
which proves all of \eqref{geo:approximation-limits}.
\end{proof}

\begin{proof}[Proof of the inequality in Theorem~\ref{intro:entropy}]
Apply Theorem~\ref{mat:entropy-smooth} to the normalized potentials
$V_j+\log Z_j$ of Lemma~\ref{geo:approximation}.  Their Hessians satisfy
\eqref{geo:hessian-bounds}, so
\[
  h(f_j)\geq m+\frac12\log\det\Cov(f_j).
\]
The covariance of $f$ is positive definite: zero variance in a nonzero
direction would concentrate this Lebesgue density on an affine hyperplane.
Thus \eqref{geo:approximation-limits} permits passage to the limit in both
sides and proves the inequality.
\end{proof}

